\subsection{Papel funcional do tipo óleo no produto final}

Cada óleo ou gordura possui um perfil próprio de ácidos graxos, e é esse perfil que define o comportamento do sabão resultante. Atributos como dureza da barra, volume e textura da espuma, capacidade de limpeza e suavidade sobre a pele variam, portanto, conforme a matéria-prima escolhida. Esse vínculo entre composição e desempenho permite tratar cada óleo como um componente de função conhecida na formulação, abrindo espaço para o uso de misturas em que as qualidades de uma matéria-prima compensam as limitações de outra \cite{araujo2015valorizacao}.

\subsubsection{Óleo de coco}

O óleo de coco possui alto teor de triglicerídeos de cadeia média saturados, com predominância de ácido láurico (C12:0) e ácido mirístico (C14:0). Em razão da saturação de seus ácidos graxos, as moléculas dos sais formados se empacotam de maneira compacta, resultando em um sabão mais duro. Por possuírem cadeias de comprimento médio, esses ácidos graxos conferem ao óleo um índice de saponificação elevado e ao produto final um alto poder de limpeza e abundante formação de espuma. Atua, portanto, como o principal agente espumante e de limpeza da formulação, ainda que, em proporções elevadas, possa ressecar a pele.

\subsubsection{Óleo de soja}

O óleo de soja resulta em ácidos graxos de cadeia longa com alto nível de insaturações, com destaque para os ácidos linoleico (C18:2) e oleico (C18:1). Como as cadeias insaturadas dificultam o empacotamento molecular, formam-se sabões mais macios, com menor poder de limpeza em comparação ao óleo de coco e teor de espuma moderado. Além disso, a presença de insaturações torna esse óleo mais suscetível à rancificação, ou seja, à oxidação dos ácidos graxos, processo que pode gerar mau odor e manchas no sabão.

\subsubsection{Óleo de rícino}

Os ácidos graxos do óleo de rícino apresentam uma característica distinta dos demais: cadeias longas com insaturação e um grupo hidroxila, com elevada composição de ácido ricinoleico (C18:1-OH). O sabão gerado a partir desse óleo possui maior polaridade e tende a aumentar a viscosidade da formulação, o que lhe confere alta solubilidade em água e uma espuma cremosa. Por conta do comportamento reológico da mistura, sua reação de saponificação pode parecer mais lenta em algumas formulações artesanais, ainda que isso dependa das condições do processo. Essas características tornam o óleo de rícino especialmente útil na produção de sabão líquido, geralmente com KOH como base.

\subsubsection{Óleo de coco babaçu}
\label{subsubsec:oleo_babacu}

O óleo de babaçu apresenta perfil rico em ácidos graxos saturados de cadeia média, com predominância de ácido láurico (C12:0) e contribuições de mirístico (C14:0) e palmítico (C16:0). Por ter composição próxima à dos óleos láuricos, como o óleo de coco, tende a formar sabões mais firmes, com boa compactação molecular e alta estabilidade da barra. Seu índice de saponificação é relativamente elevado, o que favorece uma reação eficiente e um bom rendimento de sabão em condições de bancada. Além disso, costuma gerar espuma abundante e limpeza intensa, constituindo uma alternativa interessante ao óleo de coco tradicional, embora, em proporções muito altas, também possa aumentar o potencial de ressecamento da pele.

\subsubsection{Óleo residual de cozinha}

O óleo residual de cozinha é empregado com fins de sustentabilidade e circularidade do processo produtivo. Sua composição típica é uma mistura variável de ácidos graxos livres e ésteres, decorrente da degradação sofrida durante o uso. Apesar de constituir uma alternativa econômica e ambientalmente viável, a presença de impurezas pode alterar a cor e o odor do produto final, exigindo um pré-tratamento, como a filtração, antes de sua utilização.